\section{Results}

We evaluate our encoder-based SAE steering method on bias reduction and model utility, comparing it against baseline (no steering) and direct activation steering. Our primary finding is that encoder-based SAE steering achieves statistically significant bias reduction on CrowS-Pairs while maintaining performance on StereoSet, with a modest tradeoff in language modeling utility.

\subsection{Bias Reduction Results}

\subsubsection{CrowS-Pairs Evaluation}

Table~\ref{tab:bias_results} presents anti-stereotype accuracy on the full CrowS-Pairs test set ($n=1508$). The baseline model (no steering, $\alpha=0$) achieves an anti-stereotype accuracy of $36.8\%$, indicating a preference for stereotypical completions. Our encoder-based SAE steering method with $\alpha=2.0$ improves this to $40.58\%$, representing a $3.78$ percentage point absolute improvement and a $10.3\%$ relative improvement.

We perform a one-tailed z-test to assess statistical significance, comparing the proportion of anti-stereotypical preferences between baseline and steered models. The improvement yields $z \approx 2.14$ with $p < 0.05$, confirming that the bias reduction is statistically significant.

In contrast, direct activation steering (which subtracts a mean difference vector directly from hidden states) fails to achieve comparable improvements. Across multiple $\alpha$ values tested ($-2.0$ to $2.0$), direct activation steering shows minimal variation in CrowS-Pairs performance, with anti-stereotype accuracy remaining near baseline levels ($36.63\%$ to $36.88\%$). This demonstrates that the interpretability-guided approach of using SAE encoder directions is more effective than naive vector subtraction.

\subsubsection{StereoSet Evaluation}

On the full StereoSet intrasentence validation set ($n=2106$), our encoder-based SAE steering maintains performance comparable to baseline. The baseline model achieves $60.3\%$ anti-stereotype accuracy, while encoder steering achieves $60.02\%$—a negligible change of $-0.28$ percentage points. This stability is notable given the significant improvement on CrowS-Pairs, suggesting that our method selectively reduces bias without degrading performance on this complementary benchmark.

Direct activation steering similarly shows minimal variation on StereoSet, with anti-stereotype accuracy ranging from $60.0\%$ to $60.62\%$ across different $\alpha$ values, indicating that this benchmark may be less sensitive to the steering interventions we tested.

\subsection{Utility Evaluation}

To assess whether bias reduction comes at the cost of general language modeling capability, we measure perplexity on WikiText-2. Table~\ref{tab:utility_results} shows that the baseline model achieves a perplexity of $5.697$ on a 50,000-token subset of the test set. With encoder-based SAE steering ($\alpha=2.0$), perplexity increases to $6.351$, representing an $11.5\%$ relative increase.

This utility degradation is expected, as steering interventions modify internal representations to reduce bias, which may slightly disrupt the model's learned statistical patterns. However, the increase is modest relative to the bias reduction achieved, suggesting that encoder-based SAE steering provides a favorable tradeoff between bias mitigation and model utility.

\subsection{Comparison: Encoder vs. Decoder Directions}

A key contribution of our work is demonstrating that encoder directions ($\mathbf{E}$) are more effective for bias mitigation than decoder directions ($\mathbf{D}$). While decoder directions are optimized for reconstruction, encoder directions directly capture the feature space in which bias manifests. Our encoder-based approach achieves the $10.3\%$ relative improvement on CrowS-Pairs, whereas decoder-based steering (not shown in detail) fails to achieve comparable gains. This finding highlights the importance of choosing the appropriate SAE component for interpretability-guided interventions.

\subsection{Summary}

Our results demonstrate that encoder-based SAE steering successfully reduces bias on CrowS-Pairs with statistical significance, while maintaining performance on StereoSet. The method incurs a modest utility cost ($11.5\%$ perplexity increase), which represents a reasonable tradeoff for bias mitigation. Critically, our approach outperforms direct activation steering, validating the value of interpretability-guided interventions using SAE feature decomposition.

\begin{table}[h]
\centering
\caption{Bias reduction results on CrowS-Pairs and StereoSet benchmarks. Encoder-based SAE steering ($\alpha=2.0$) achieves significant improvement on CrowS-Pairs while maintaining performance on StereoSet.}
\label{tab:bias_results}
\begin{tabular}{lcc}
\toprule
Method & CrowS-Pairs & StereoSet \\
 & Anti-Stereotype Acc. & Anti-Stereotype Acc. \\
\midrule
Baseline ($\alpha=0$) & $36.8\%$ & $60.3\%$ \\
Encoder SAE Steering ($\alpha=2.0$) & $\mathbf{40.58\%}$ & $60.02\%$ \\
Direct Activation Steering ($\alpha=2.0$) & $36.63\%$ & $60.50\%$ \\
\bottomrule
\end{tabular}
\end{table}

\begin{table}[h]
\centering
\caption{Utility evaluation on WikiText-2 perplexity. Encoder-based SAE steering incurs a modest $11.5\%$ relative increase in perplexity.}
\label{tab:utility_results}
\begin{tabular}{lc}
\toprule
Method & WikiText-2 PPL \\
\midrule
Baseline ($\alpha=0$) & $5.697$ \\
Encoder SAE Steering ($\alpha=2.0$) & $6.351$ \\
\bottomrule
\end{tabular}
\end{table}

